\chapter{一页式学习回路设计表}

这张表用于重做一个 45--90 分钟教学单元。先填写“学生要做什么”，最后才填写“教师讲什么”。

\begin{description}[style=nextline,leftmargin=3.2cm]
  \item[值得求解的问题] 学生为什么愿意解决？问题有哪些真实约束？
  \item[目标判断] 课后学生应能作出什么专业判断，而不仅是记住什么？
  \item[首次作品] 15 分钟内能产出什么可见结果？允许粗糙，但必须可检查。
  \item[反馈来源] 结构响应、同伴、量规、教师或 AI 分别反馈什么？
  \item[第二次修改] 学生必须根据反馈改变哪一处，并保存前后版本。
  \item[解释与答辩] 学生怎样说明“改了什么、依据是什么、仍不确定什么”？
  \item[AI 边界] AI 可以生成哪些材料？哪些判断必须由学生亲自完成？
  \item[验收证据] 初稿、反馈、修改稿、来源和口头解释是否齐全？
\end{description}

\section*{最小示例：简支梁跨中弯矩}

不要先让学生抄写公式。先给三座跨度不同、截面相同的简支梁，让学生预测哪一座最危险并画出弯矩图形状；再用计算或交互工具反馈；最后要求学生用 $qL^2/8$ 解释跨度为何如此敏感。AI 可以帮助画图和检查代数，但学生必须说明边界条件、单位和图形特征。

\section*{验收标准}

一个合格单元至少包含一次真实修改。如果学生只生成一次答案、教师只给一次分数，这不是学习回路，只是一次提交。
